\section{The small-field remainder of the pure nonlinear push}
\label{sec:pure-action-small-field}
We combine the truncated stationary integral with the full positive
push and the classical matching conditions. The full logarithm and
the stationary logarithm are kept distinct until their difference
has been estimated.

\subsection{From an exterior bound to a logarithmic error}
Let $Z(x)$ be the actual pure nonlinear push in coarse logarithmic
coordinates, and let $\Omega>0$ be its parameter-independent
normalizing factor. The stationary construction supplies a real
normalizing exponent $\Phi(x)$ and the positive truncated integral
$R(\beta^{-1/2},x)$ of
Section~\ref{sec:truncated-stationary-log}. Suppose its exterior
estimate is
\begin{equation}\label{eq:pure-exterior-receipt}
 0\le Z(x)-\Omega e^{-\Phi(x)}R(\beta^{-1/2},x)
                          \le\Omega e^{-\Phi(x)}\eta.
\end{equation}
The integral over the retained stationary region is exactly the
middle term, with the same moving Hessian and Haar-chart factor.
The complementary integral is nonnegative. Its upper bound in
\eqref{eq:pure-exterior-receipt} is a substantive exterior estimate,
not a consequence of positivity.

Assume $R\ge a_0>0$. Then $Z>0$ and, for
\[
 E(x)=\log Z(x)-\log\Omega+\Phi(x)
                              -\log R(\beta^{-1/2},x),
\]
we have
\begin{equation}\label{eq:pure-log-value}
                             0\le E(x)\le\eta/a_0.
\end{equation}
Indeed the ratio $Z/(\Omega e^{-\Phi}R)$ lies between $1$ and
$1+\eta/R$, and $0\le\log(1+s)\le s$ for $s\ge0$.
In the construction $a_0=e^{-1}/2$ and
$\eta=\exp(-cP_F^2)$ with a fixed $c>0$.

\subsection{Polynomial derivative bounds for the full logarithm}
The coarse field enters the positive integral through the smoothing
kernel. At a fixed fine configuration and fixed tag its contribution
has the form
\[
                 U(x,\omega)=a(\omega)e^{p(x,\omega)},
                       \qquad a(\omega)\ge0.
\]
The fine density and tag likelihood are held in $a$. The exponent
$p$ is minus $\beta$ times the finite sum of kernel penalties
between the varied coarse bonds and the fixed predicted bonds.
On a fixed logarithmic ball, smoothness of the group exponential
and the penalty gives a common derivative bound
\[
              \|p(\cdot,\omega)\|_{J^k,x}
                              \le C_k|\beta|\,n_b,
\]
where $n_b$ bounds the number of coarse bonds. It is uniform in the
held fine field because the group representation is unitary and
that field enters only through fixed group multiplication.

Subtract $p(x,\omega)$ before differentiating the exponential.
The finite partition formula gives
\[
               \|U(\cdot,\omega)\|_{J^k,x}
                                      \le K_k U(x,\omega),
\]
where $K_k$ is a fixed polynomial in $1+|\beta|n_b$.
Finite tag summation preserves this relative bound. Differentiation
under the Haar integral is justified by the bounded fine density
and kernel on the compact configuration space, together with the
uniform derivative bound on a neighborhood. Consequently
\begin{equation}\label{eq:pure-relative-derivatives}
                            \|Z\|_{J^k,x}\le K_k Z(x).
\end{equation}

At a point where $Z(x)>0$, normalize by that one value:
$f(y)=Z(y)/Z(x)$. Then $f(x)=1$ and its jets are bounded by $K_k$.
The logarithm chain rule at one bounds the jets of
$\log Z(y)-\log Z(x)$ by a fixed polynomial in $K_k$. This argument
never divides an absolute derivative bound by an uncontrolled small
partition function; the factor $Z(x)$ has already canceled in
\eqref{eq:pure-relative-derivatives}.

The other terms defining $E$ have polynomial derivative bounds:
$\Phi$ consists of the stationary action multiplied by $\beta$,
the stationary extra exponent and the relative log determinant;
the truncated logarithm has the bounds proved in the preceding
section. For auxiliary volumes polynomial in the block radius and
$\beta$ in its allowed band, these give a bound $z^d$ through any
fixed positive order, for some fixed $d$, where $z=\beta_F(x_{\rm scale})$.
The order-zero bound for $E$ follows directly from
\eqref{eq:pure-log-value}.
The exponentially small value bound
\eqref{eq:pure-log-value} holds on a microscopic ball of radius
$z^{-1}$ about each tested coarse point. Applying the interpolation
lemma of the preceding section yields, for every prescribed $e$,
\begin{equation}\label{eq:pure-log-jets}
                              \|E\|_{J^j}\le z^e.
\end{equation}
This uses the exterior estimate throughout that ball. An estimate
only at its center would not suffice.

\subsection{The exact decomposition and the two cancellations}
Write $a_{\rm fine}(x)$ for the fine action evaluated on the actual
stationary branch, $a_{\rm coarse}(x)$ for the coarse classical
action, and $d_{\rm coarse}(x)$ for the prescribed coarse quantum
and reset contribution. Define
\[
 C(x)=a_{\rm coarse}(x)-a_{\rm fine}(x),\qquad
 W(x)=d_{\rm coarse}(x)
             -\psi(x,\rho(x))-\tfrac12\log\det M(x).
\]
Thus the normalizing exponent is
$\Phi=\beta a_{\rm fine}+\psi(x,\rho(x))+\tfrac12\log\det M(x)$.
The actual corrected logarithm
\[
                  G(x)=\log Z(x)+\beta a_{\rm coarse}(x)
                                      +d_{\rm coarse}(x)
\]
satisfies the exact identity
\begin{equation}\label{eq:pure-action-decomposition}
                G=\log\Omega+L+E+\beta C+W,
                 \qquad L(x)=\log R(\beta^{-1/2},x).
\end{equation}
Every function in this identity uses the same stationary branch,
coefficient data and finite torus.

Classical matching through degree four says
$D^rC(0)=0$ for $0\le r\le4$. Suppose its derivatives of orders
$1$ through $n\ge5$ on the ball of radius $h\le1$ are bounded
by $K_C$. Taylor's formula, applied also to each derivative of
order less than five, gives
\begin{equation}\label{eq:pure-classical-fifth}
                    C\in\mathcal B_n(K_Ch^5,h^{-1}).
\end{equation}
For orders at least five the original derivative bound suffices,
because $h^r\le h^5$; for lower orders integrate the fifth derivative
along the radial segment, using the zero jets at its starting point.

There is a distinct cancellation for the linear part of $W$.
The actual corrected logarithm $G$ is invariant under the prescribed
lattice translations and signed coordinate symmetries. For each
coordinate basis vector supported on the bond $(b,\mu)$, central
inversion followed by translation by $2b+e_\mu$ fixes that
unoriented bond and reverses its orientation. It therefore sends
the single-coordinate logarithmic vector $e_i$ to $-e_i$. Differentiating invariance at zero gives
$DG(0)e_i=-DG(0)e_i$, hence $DG(0)=0$.
Since $DC(0)=0$, \eqref{eq:pure-action-decomposition} implies
\[
                         DW(0)=-DL(0)-DE(0).
\]
In particular, if the relevant jets of $L,E$ are bounded by $a,v$,
then $\|DW(0)\|\le a+v$. This cancellation uses the symmetry of
the full logarithm, rather than asserting that its separate
stationary terms are individually even in every coarse coordinate.

The finite-volume quadratic matching estimate supplies a further
bound $\|D^2W(0)\|\le q$. Its operator norm includes the conversion
from the tested momentum waves to the actual coordinate space.
If the number of coordinates is at most $M$, that conversion has
a fixed polynomial cost in $M$. Thus the exponentially small
finite-volume bound has the form
\[
                 q\le P(M)\Theta\Lambda e^{-\nu\ell},
                       \qquad \Theta,\nu>0,
\]
where $\ell$ is the auxiliary torus side and $P$ is a fixed
nonnegative polynomial including the coordinate norm conversion.
This quadratic matching is a separate input from classical
matching through degree four.

If derivatives of $W$ of orders three through $n$ are bounded by
$K_W$ on the same ball, Taylor's formula now gives
\begin{equation}\label{eq:pure-quantum-third}
 W-W(0)\in\mathcal B_n
       \bigl(K_Wh^3+3\{(a+v)h+qh^2\},h^{-1}\bigr).
\end{equation}
For orders zero, one and two, integrate the third derivative and
retain the displayed first and second jets. For higher orders,
$h^r\le h^3$ pays the original bound. The factor three is a
common upper bound for all these finite Taylor coefficients.

\subsection{The final small-field price}
Take $h=z^{-7/16}$ and assume $\beta\in[z/2,2z]$.
The stationary estimate gives
$a=C(x_{\rm scale})\beta^{-1}+z^{-2}$ with $C$ of upper exponent
zero; \eqref{eq:pure-log-jets} gives $v=z^{-2}$. On the radial
segment to the tested point, the centered stationary logarithm has
price $ah$ at scale $h^{-1}$, by integrating its first derivative;
the higher derivative bounds give the same price since $h\le1$.
The centered error has price $2v$.
Using \eqref{eq:pure-action-decomposition}--\eqref{eq:pure-quantum-third}
therefore gives
\begin{equation}\label{eq:pure-remainder-budget}
 G-G(0)\in\mathcal B_j
 \bigl(ah+2v+\beta K_Ch^5
                +K_Wh^3+3\{(a+v)h+qh^2\},h^{-1}\bigr)
                       \qquad(1\le j\le n).
\end{equation}
All constants refer to the same original data.

The coefficients $K_C,K_W$, the coordinate bound $M$, and the
loop price $\Lambda$ have upper exponent zero; $\Theta,\nu$ are
fixed positive constants. The auxiliary
side is at least the prescribed logarithmic-square scale, so its
exponential finite-volume factor pays every fixed inverse power;
in particular $q$ may be bounded by a function of upper exponent
$-1$. The terms in the displayed amplitude have upper exponents
at most
\[
          -23/16,\quad -2,\quad -19/16,\quad
                         -21/16,\quad -23/16,\quad -15/8.
\]
Their maximum is $-19/16<-23/20$. After one common increase of the
stage threshold, the amplitude is at most $z^{-23/20}$.
The estimate holds at every healthy coarse point with the two
microscopic margins needed for the comparison ball and the radial
segment. It supplies the centered chart estimate used in
Section~\ref{sec:healthy-boundary-jets}. Its remaining hypotheses are
the actual stationary-branch bounds, classical and quadratic
matching, and the exterior estimate
\eqref{eq:pure-exterior-receipt}; none is replaced by a comparison
with a different finite-volume model.
